% tube-algorithm.tex — standalone section, to be \input'd from algorithms.tex
% Conventions: thesis/CLAUDE.md
%
% Sources: Jörn's dictation (tube-notes.md, 2026-02-18) + CH2021 (rotation number,
%          symplectic flow graph, transition matrices) + HK2017 (simple orbit theorem).
%          All content is agent-written and unreviewed unless marked.
%
% Structure:
%   \subsection{Setting} — symplectic polytopes, directed 2-faces
%   \subsection{Tubes} — definition, convexity, data structure
%   \subsection{Extension} — step maps, tube extension formula
%   \subsection{Rotation Number} — universal cover, additivity, increments
%   \subsection{Algorithm} — pruning lemmas, closing step, algorithm box
%   \subsection{Correctness} — sketch, Type 2 remark

% [TODO: JÖRN - This section is agent-written from your dictation (tube-notes.md).
%   Review the full section for mathematical accuracy before citing it as thesis content.
%   Key items needing your check are marked with [TODO: JÖRN - ...] below.]

\section{Tube Algorithm}\label{sec:tube-algorithm}

% [TODO: JÖRN - intro paragraph once the section is stable.]

The \emph{tube algorithm} is an alternative to Algorithm~\ref{alg:ehz}
for computing $c_{\mathrm{EHZ}}(K)$ on a class of polytopes
we call \emph{symplectic polytopes} (Definition~\ref{def:symplectic-polytope}).
The idea is to search for the minimum-action orbit iteratively:
instead of enumerating all combinatorial types up front,
we build up families of trajectories (``tubes'') one facet at a time,
pruning branches that cannot contain a minimum-action orbit.

% -------------------------------------------------------------------------
\subsection{Setting}\label{subsec:tube-setting}
% -------------------------------------------------------------------------

We recall the face structure of the polytopes we work with
and introduce the directed flow structure that the algorithm exploits.

\begin{definition}[Faces of a polytope]\label{def:faces}
% QC: face = standard convex geometry definition (extremal subset).
%   Equivalently: F = K ∩ H for a supporting hyperplane H, but different H can give the same F.
%   k-face = affine dimension k. Closed set, open relative interior.
%   For a polytope in irredundant H-rep, faces at every dimension 0-3 exist.
%   Facets (3-faces) are F_i per Definition~\ref{def:facets}.
%   2-faces: by Lemma~\ref{lem:2-face-structure}, exactly one pair {i,j} per 2-face.
%   1-faces and 0-faces: may lie in 3 or more facets (no uniqueness, see Remark~\ref{rem:lower-faces}).
% Downstream: face data needed to enumerate tube transitions: only 2-faces matter (Type 1).
Let $K \subset \mathbb{R}^4$ be a polytope
with irredundant $H$-representation
$\{(n_i, h_i)\}_{i=1}^{F}$ (Definition~\ref{def:facets}).
A \emph{face} of $K$ is a convex subset $G \subset K$
such that for any line segment $[x, y] \subset K$
whose relative interior meets $G$,
the entire segment lies in $G$.
A \emph{$k$-face} is a face of affine dimension $k$.
For polytopes, every face is an intersection of facets:
\begin{itemize}
\item \emph{3-faces} (facets) are $F_1, \ldots, F_F$
  (Definition~\ref{def:facets}).
\item A \emph{2-face} is an intersection $F_i \cap F_j$
  that is $2$-dimensional.
  By Lemma~\ref{lem:2-face-structure},
  each $2$-face arises from exactly one pair $\{i,j\}$;
  we write $F_{ij}$ for this face.
\item A \emph{1-face} (edge) is an intersection
  $F_{i_1} \cap \cdots \cap F_{i_m}$ that is $1$-dimensional,
  for some $m \geq 3$.
  A \emph{0-face} (vertex) is such an intersection
  that is $0$-dimensional, for some $m \geq 4$.
  Unlike $2$-faces, there is no canonical minimal
  set of generating facets
  (Remark~\ref{rem:lower-faces}).
\end{itemize}
Every face is closed; $\operatorname{int}(F)$ denotes its relative interior.
\end{definition}

\begin{lemma}[2-face uniqueness]\label{lem:2-face-structure}
% QC: every 2-face is contained in exactly 2 facets. Proof: normal cone argument.
%   Normal cone at int(F_{ij}) lives in a 2D normal space; in 2D, a proper
%   pointed cone has exactly 2 extreme rays; irredundancy prevents additional facets.
%   For 1-faces: normal space is 3D, so normal cone can have ≥3 extreme rays.
% Downstream: F_{ij} is unambiguously labeled by the pair {i,j}. No 2-face is F_i ∩ F_j ∩ F_k.
Every $2$-face of $K$ is contained in exactly two facets.
\end{lemma}

\begin{proof}
Let $F$ be a $2$-face and let $x \in \operatorname{int}(F)$.
The outward normal cone
$N_x K \coloneqq \{ u \in \mathbb{R}^4 : \langle u, y - x \rangle \leq 0 \;\forall\, y \in K \}$
is a closed convex cone.
Since $\dim F = 2$ and $\dim K = 4$, the normal cone lives in the
$2$-dimensional linear subspace $(\operatorname{aff} F)^\perp$.
A proper pointed cone in $\mathbb{R}^2$ is bounded by
exactly two extreme rays.
In the irredundant $H$-representation,
each ray corresponds to exactly one facet:
a facet $F_i$ contains $x$ if and only if $n_i$ is in the cone
(since $\langle x, n_i \rangle = h_i$ iff $F_i \ni x$).
Irredundancy prevents any $n_i$ from lying strictly
in the interior of the cone generated by the other normals
(such a facet would be implied by the others).
Hence the cone has exactly two extreme rays,
corresponding to exactly two facets containing $F$.
\end{proof}

\begin{remark}[Lower-dimensional faces]\label{rem:lower-faces}
% QC: 1-faces and 0-faces are not uniquely determined by a pair/triple of facets.
%   A 1-face lies in 3 or more facets: its normal cone lives in R^3 and
%   can have many extreme rays. Similarly 0-faces in R^4.
%   These correspond to Type 2 orbits (Remark~\ref{rem:type-2}).
Unlike $2$-faces, a $1$-face or $0$-face may lie in three or more facets.
The dimension argument of Lemma~\ref{lem:2-face-structure}
does not apply: the normal cone of a $1$-face lives in $\mathbb{R}^3$
and can have arbitrarily many extreme rays.
\end{remark}

\begin{definition}[Symplectic polytope]\label{def:symplectic-polytope}
% QC: each 2-face is either symplectic or Lagrangian as a 2D subspace.
%   Symplectic: omega_0 non-degenerate on TF_{ij}. Lagrangian: omega_0|_{TF_{ij}}=0.
%   omega_0|_{TF_{ij}}=0 iff omega_0(n_i,n_j)=0 (since TF_{ij} ⊂ n_i^perp ∩ n_j^perp,
%   and n_i^perp ∩ n_j^perp is a 2D symplectic subspace iff omega_0(n_i,n_j) != 0).
%   The name "symplectic polytope" is from \cite{CH2021}, Definition~1.x.
% Downstream: check omega_0(n_i, n_j) != 0 for all 2-dimensional F_i ∩ F_j.
%   Skip pairs where F_i ∩ F_j is lower-dimensional (empty, 0D, or 1D).
A \emph{Lagrangian $2$-face} is a $2$-face $F_{ij}$ on which
$\omega_0|_{T F_{ij}} = 0$, equivalently $\omega_0(n_i, n_j) = 0$.
% [TODO: JÖRN - verify the equivalence: TF_{ij} subset n_i^perp ∩ n_j^perp,
%   and ω_0 restricted to this 2D space is zero iff ω_0(n_i,n_j) = 0.]
A \emph{symplectic polytope} is a polytope $K \subset \mathbb{R}^4$
that has no Lagrangian $2$-face:
\[
  \omega_0(n_i, n_j) \neq 0
  \qquad \text{for all pairs } (i,j)
  \text{ such that } F_i \cap F_j \text{ is a $2$-face.}
\]
\end{definition}

\begin{definition}[Directed $2$-face]\label{def:directed-2-face}
% QC: sign of omega_0(n_i,n_j) determines flow direction through int(F_{ij}).
%   From CH2021 Lemma~\ref{lem:EinEout} (using i=quaternionic J_0): if omega_0(n_i,n_j)>0
%   then R_i exits F_i through F_{ij}, i.e. trajectories on F_i leave through F_{ij} into F_j.
%   Equivalently: a trajectory with Reeb velocity R_i, reaching int(F_{ij}), continues into F_j.
% Downstream: directed edge i->j exists iff omega_0(n_i,n_j) > 0 AND F_i ∩ F_j is 2-dimensional.
%   Precompute this adjacency relation for the algorithm.
Let $K$ be a symplectic polytope and $F_{ij}$ a $2$-face.
Since $\omega_0(n_i, n_j) \neq 0$, we define the \emph{direction}:
\[ F_i \to F_j \iff \omega_0(n_i, n_j) > 0. \]
By \cite{CH2021}, Lemma~{lem:EinEout}, a Reeb trajectory with velocity $R_i$
exits the interior of $F_i$ through $\operatorname{int}(F_{ij})$
if and only if $F_i \to F_j$.
Equivalently, the Reeb flow at $F_i$ transitions to $F_j$
at any breakpoint in $\operatorname{int}(F_{ij})$.
\end{definition}

% -------------------------------------------------------------------------
\subsection{Tubes}\label{subsec:tubes}
% -------------------------------------------------------------------------

We now define the objects the algorithm maintains:
sets of pure Reeb trajectories sharing a facet-visit sequence.

\begin{definition}[Pure Reeb trajectory]\label{def:pure-reeb-trajectory}
% QC: a special case of Definition~\ref{def:reeb-orbit-polytope}: velocity is a pure Reeb
%   vector R_{sigma(i)}, not a convex mixture, on each interval (t_{i-1},t_i).
%   gamma(t) in F_{sigma(i)} for t in [t_{i-1}, t_i] (by R_i tangent to bdry(K) along F_i).
%   t_1 = 0 is a normalisation (free to choose by translation).
% Downstream: trajectory determined by starting point gamma(0) in F_{sigma(1)} ∩ F_{sigma(2)}.
%   All other positions computable by integrating R_{sigma(i)} forward.
Let $\sigma = (\sigma(1), \ldots, \sigma(k))$ be a sequence of facet indices with $k \geq 2$.
A \emph{pure Reeb trajectory of type $\sigma$} is a curve
$\gamma \colon [t_0, t_k] \to \partial K$ with
$t_0 < t_1 < \cdots < t_k$, $t_1 = 0$, and
\[
  \dot{\gamma}(t) = R_{\sigma(i)}
  \qquad \text{for } t \in (t_{i-1}, t_i),
  \quad i = 1, \ldots, k.
\]
The times $t_2, \ldots, t_k$ are determined by the trajectory:
$t_i$ is the first time after $t_{i-1}$ that $\gamma$ exits $F_{\sigma(i)}$.
\end{definition}

\begin{definition}[Tube]\label{def:tube}
% QC: the tube T(sigma) is a convex family (Lemma~\ref{lem:tube-convex}).
%   T(sigma) may be empty (no valid trajectory exists).
%   The tube is parameterized by its starting point gamma(0) ∈ Start(sigma) ⊂ F_{sigma(1)} ∩ F_{sigma(2)}.
The \emph{tube} $T(\sigma)$ of a type $\sigma = (\sigma(1), \ldots, \sigma(k))$
is the set of all pure Reeb trajectories of type $\sigma$:
\[
  T(\sigma) \coloneqq \bigl\{
    \gamma \colon [t_0, t_k] \to \partial K
    : \gamma \text{ is a pure Reeb trajectory of type } \sigma
  \bigr\}.
\]
\end{definition}

\begin{lemma}[Convexity of a tube]\label{lem:tube-convex}
% QC: the key is that once we fix sigma, the dynamics are globally affine:
%   gamma_lambda = lambda * gamma_1 + (1-lambda) * gamma_2 has velocity
%   lambda R_i + (1-lambda) R_i = R_i on each interval. Valid: velocities match,
%   and affine combination stays on boundary because F_i is affine and gamma_1,gamma_2 in F_i.
% [TODO: JÖRN - the argument uses that affine combinations of points on F_i stay on F_i;
%   F_i is closed and convex so lambda x + (1-lambda)y in F_i if x,y in F_i. Correct.]
If $\gamma_1, \gamma_2 \in T(\sigma)$ and $\lambda \in [0,1]$, then
$\lambda \gamma_1 + (1-\lambda)\gamma_2 \in T(\sigma)$.
\end{lemma}

\begin{proof}
On each interval $(t_{i-1}(\gamma_1), t_i(\gamma_1))$,
the velocity $\dot{\gamma}_1(t) = R_{\sigma(i)}$ is constant,
and similarly for $\gamma_2$.
Since $F_{\sigma(i)}$ is a convex set,
the affine combination $\gamma_\lambda(t) \coloneqq \lambda\gamma_1(t) + (1-\lambda)\gamma_2(t)$
remains in $F_{\sigma(i)}$ wherever $\gamma_1$ and $\gamma_2$ do.
The velocity $\dot{\gamma}_\lambda(t) = R_{\sigma(i)}$ on each interval.
The breakpoints may shift, but $\gamma_\lambda$ is still a pure Reeb trajectory of type $\sigma$.
Hence $\gamma_\lambda \in T(\sigma)$.
\end{proof}

\begin{definition}[Start and End of a tube]\label{def:tube-endpoints}
% QC: Start = gamma(t_1 = 0) in F_{sigma(1)} ∩ F_{sigma(2)}: the 2-face entered at t_1=0.
%   End = gamma(t_{k-1}) in F_{sigma(k-1)} ∩ F_{sigma(k)}: the LAST INNER breakpoint.
%   Note: t_0 and t_k (= outer breakpoints) are gamma-dependent; t_1,...,t_{k-1} are
%   the inner breakpoints. Start and End are convex sets (by tube convexity).
%   The map gamma <-> gamma(0) is a bijection: gamma is determined by its start.
% Downstream: Start and End are convex polygons in 2D local coordinates of F_{sigma(1)} ∩ F_{sigma(2)}
%   (resp. F_{sigma(k-1)} ∩ F_{sigma(k)}). Stored as vertex lists.
For a tube $T(\sigma)$ with $k \geq 3$:
\begin{align*}
  \operatorname{Start}(\sigma) &\coloneqq
    \{ \gamma(t_1) : \gamma \in T(\sigma) \}
    \subset F_{\sigma(1)} \cap F_{\sigma(2)}, \\
  \operatorname{End}(\sigma) &\coloneqq
    \{ \gamma(t_{k-1}) : \gamma \in T(\sigma) \}
    \subset F_{\sigma(k-1)} \cap F_{\sigma(k)}.
\end{align*}
Both sets are closed convex polygons,
by Lemma~\ref{lem:tube-convex} and compactness of $K$.
\end{definition}

\begin{definition}[Tube data structure]\label{def:tube-data}
% QC: (Start, End, phi, a, rho) encode all information needed for extension and closing.
%   phi is determined by gamma(0) |-> gamma(t_{k-1}): an affine bijection Start -> End
%   (globally affine since each step is affine and composition of affines is affine).
%   a(y) = total time from t_1 to t_{k-1} for the trajectory ending at y in End
%         = action of gamma|_{[t_1, t_{k-1}]} (since A(gamma) = T by contact normalization).
%   rho = rotation number accumulated from t_1 to t_{k-1} (defined in §Rotation).
% Downstream: phi stored as 2x2 matrix + offset in local coords.
%   a stored as affine function (gradient + constant) in local coords of End.
%   rho stored as a single real number.
A tube $T(\sigma)$ is represented by the following data:
\begin{enumerate}
\item $\operatorname{Start} \subset F_{\sigma(1)} \cap F_{\sigma(2)}$:
  closed convex polygon (in local $2$-dimensional coordinates of the $2$-face).
\item $\operatorname{End} \subset F_{\sigma(k-1)} \cap F_{\sigma(k)}$:
  closed convex polygon (in local $2$-dimensional coordinates).
\item $\varphi \colon \operatorname{Start} \to \operatorname{End}$:
  the affine map sending $\gamma(0)$ to $\gamma(t_{k-1})$.
\item $a \colon \operatorname{End} \to \mathbb{R}$:
  the affine function $a(y) = t_{k-1}$,
  the total elapsed time (= action) of $\gamma|_{[t_1, t_{k-1}]}$
  for the trajectory with $\gamma(t_{k-1}) = y$.
\item $\rho \in \mathbb{R}$:
  the rotation number accumulated on $(t_1, t_{k-1})$;
  see Definition~\ref{def:rotation-number} and Remark~\ref{rem:rotation-polytope}.
\end{enumerate}
The map $\varphi$ is affine (composition of affine step maps; see §\ref{subsec:tube-extension})
and $a$ is affine in local coordinates.
\end{definition}

\begin{remark}[Initialization for $k = 2$]\label{rem:tube-init}
% QC: for k=2, tube = all trajectories entering at F_{sigma(1)} ∩ F_{sigma(2)}.
%   Start = End = F_{sigma(1)} ∩ F_{sigma(2)} (the full 2-face, no constraint yet).
%   phi = id, a = 0 (no time elapsed between t_1 and t_{k-1} = t_1), rho = 0.
%   Extension to k=3 is the first step that produces nontrivial data.
For a tube of length $k = 2$ (type $(\sigma(1), \sigma(2))$),
we set $\operatorname{Start} = \operatorname{End} = F_{\sigma(1)} \cap F_{\sigma(2)}$,
$\varphi = \operatorname{id}$, $a \equiv 0$, and $\rho = 0$.
\end{remark}

% -------------------------------------------------------------------------
\subsection{Extension}\label{subsec:tube-extension}
% -------------------------------------------------------------------------

We now describe how to extend a tube by one facet
and derive the update formulas for the data structure.

\begin{definition}[Step map for a triple]\label{def:step-data}
% QC: Phi_{ijl} maps x in H_i ∩ H_j to gamma(t_k) in H_j ∩ H_l under flow by R_j.
%   Explicit formula: Phi_{ijl}(x) = x + t(x) R_j where t(x) = (h_l - <x,n_l>) / <R_j,n_l>.
%   Denominator <R_j,n_l> = (2/h_j) omega_0(n_j,n_l) > 0 (since F_j -> F_l means omega_0>0).
%   So the denominator is strictly positive; Phi is well-defined and affine.
%   Delta a_{ijl}(x) = t(x) (time elapsed = action increment).
%   Delta rho_{jl} = rotation increment at the j->l breakpoint; see §Rotation.
%   All three precomputed from (n_i,h_i), (n_j,h_j), (n_l,h_l).
% Downstream: Phi_{ijl}: R^4 -> R^4, affine. DeltaA_{ijl}: R^4 -> R, affine.
%   DeltaRho_{jl}: precomputed scalar per directed 2-face j->l.
%   Test: Phi_{ijl}(x) in H_j (always) and in H_l (by construction).
For a triple $(i,j,l)$ with $F_i \to F_j$ and $F_j \to F_l$,
define the \emph{step map}
$\Phi_{ijl} \colon H_i \cap H_j \to H_j \cap H_l$ by
\begin{equation}\label{eq:step-map}
  \Phi_{ijl}(x) \coloneqq x + t_{ijl}(x)\, R_j,
  \qquad
  t_{ijl}(x) \coloneqq \frac{h_l - \langle x, n_l \rangle}{\langle R_j, n_l \rangle},
\end{equation}
where the denominator satisfies
$\langle R_j, n_l \rangle = (2/h_j)\,\omega_0(n_j, n_l) > 0$
(since $F_j \to F_l$ means $\omega_0(n_j, n_l) > 0$).
The map $\Phi_{ijl}$ is affine and bijective.
The associated \emph{action increment} is the affine function
\begin{equation}\label{eq:action-increment}
  \Delta a_{ijl}(x) \coloneqq t_{ijl}(x),
\end{equation}
the time elapsed flowing from $x$ to $\Phi_{ijl}(x)$ along $R_j$.
The \emph{rotation increment} $\Delta\rho_{jl} \in (0, \tfrac{1}{2})$
at the $j \to l$ breakpoint is defined in Definition~\ref{def:rotation-increment}.
\end{definition}

\begin{definition}[Tube extension]\label{def:tube-extension}
% QC: given tube T(sigma(1),...,sigma(k)) and new facet sigma(k+1) with F_{sigma(k)} -> F_{sigma(k+1)},
%   the extended tube T(sigma(1),...,sigma(k+1)) has data computed by:
%   End' = Phi(End) ∩ (F_{sigma(k)} ∩ F_{sigma(k+1)})   [physical constraint]
%   phi' = Phi o phi                                       [composed affine maps]
%   Start' = phi'^{-1}(End')                              [pulled-back start]
%   a'(y') = a(phi^{-1}(Phi^{-1}(y'))) + DeltaA(Phi^{-1}(y'))  [accumulated action]
%   rho' = rho + DeltaRho_{sigma(k), sigma(k+1)}          [accumulated rotation]
%   Here Phi = Phi_{sigma(k-1), sigma(k), sigma(k+1)}.
% [TODO: JÖRN - verify the formula for Start': Start' = phi'^{-1}(End') means we
%   pull back the End constraint to the Start face. Since phi' is affine and bijective,
%   this gives a polygon in F_{sigma(1)} ∩ F_{sigma(2)}. Correct?]
% Downstream: intersection End' = Phi(End) ∩ F_{sigma(k)} ∩ F_{sigma(k+1)} is a polygon
%   intersection in 2D (Phi(End) is a polygon in H_{sigma(k)} ∩ H_{sigma(k+1)},
%   intersect with the closed face F_{sigma(k)} ∩ F_{sigma(k+1)} via H-rep constraints).
Let $T_k = T(\sigma(1), \ldots, \sigma(k))$ be a tube
and $\sigma(k+1)$ satisfy $F_{\sigma(k)} \to F_{\sigma(k+1)}$.
Write $\Phi \coloneqq \Phi_{\sigma(k-1),\sigma(k),\sigma(k+1)}$
for the step map (Definition~\ref{def:step-data}).
The extended tube $T(\sigma(1), \ldots, \sigma(k+1))$ has data:
\begin{align}
  \operatorname{End}' &\coloneqq
    \Phi(\operatorname{End}) \cap \bigl( F_{\sigma(k)} \cap F_{\sigma(k+1)} \bigr),
    \label{eq:end-prime}\\
  \varphi' &\coloneqq \Phi \circ \varphi,
    \label{eq:phi-prime}\\
  \operatorname{Start}' &\coloneqq (\varphi')^{-1}(\operatorname{End}'),
    \label{eq:start-prime}\\
  a'(y') &\coloneqq
    a\bigl(\varphi^{-1}(\Phi^{-1}(y'))\bigr)
    + \Delta a_{\sigma(k-1),\sigma(k),\sigma(k+1)}\bigl(\Phi^{-1}(y')\bigr),
    \label{eq:a-prime}\\
  \rho' &\coloneqq \rho + \Delta\rho_{\sigma(k),\sigma(k+1)}.
    \label{eq:rho-prime}
\end{align}
\end{definition}

\begin{remark}[Interpretation]\label{rem:extension-interpretation}
The intersection in~\eqref{eq:end-prime} enforces that
$\gamma(t_k) \in F_{\sigma(k)} \cap F_{\sigma(k+1)}$
is geometrically in $K$ (the image under the step map
$\Phi$ might overflow the actual face).
Equation~\eqref{eq:start-prime} propagates this constraint back to the start:
a trajectory can reach $\operatorname{End}'$ only if it started in $\operatorname{Start}'$.
The action formula~\eqref{eq:a-prime} accumulates the time
from $t_1 = 0$ to $t_k$ as a function of the new End point.
\end{remark}

% -------------------------------------------------------------------------
\subsection{Rotation Number}\label{subsec:tube-rotation}
% -------------------------------------------------------------------------

The rotation number is the key pruning tool for the search.
We follow \cite{CH2021}, \S 2.

\begin{definition}[Rotation number in the universal cover]\label{def:rotation-number}
% QC: Sp(2) = group of 2x2 symplectic matrices. Universal cover Sp~(2) is a connected
%   simply-connected Lie group. The universal cover fibres over S^1 (via the rotation
%   of the matrix in polar form). The rotation number of a path f: [0,1] -> Sp(2)
%   with f(0)=I is the R-valued lift rho(f) = rho~(f~(1)) where f~ is the lift to Sp~(2).
%   This rho is independent of the choice of path within its homotopy class rel endpoints.
%   Additivity: rho(f1 * f2) = rho(f1) + rho(f2) for composable paths in Sp~(2).
% [TODO: JÖRN - verify this is the exact same rotation number used in CH2021 §2.
%   CH2021 uses a "quaternionic trivialization" to define a canonical rotation number for
%   Reeb orbits on symplectic polytopes in R^4. Check that our description matches.]
The \emph{symplectic group} $\operatorname{Sp}(2)$ is the group of
$2 \times 2$ symplectic matrices over $\mathbb{R}$.
The \emph{universal cover} $\widetilde{\operatorname{Sp}}(2)$ is the unique
connected simply-connected Lie group covering $\operatorname{Sp}(2)$.
The covering map factors through a fibration over $S^1$
(the rotation component of the polar decomposition).
For a path $f \colon [0,1] \to \operatorname{Sp}(2)$ with $f(0) = I$,
the \emph{rotation number} is the lift
$\rho(f) \in \mathbb{R}$ of $f$ to $\widetilde{\operatorname{Sp}}(2)$,
projected to the $\mathbb{R}$-factor of the cover.
\end{definition}

\begin{remark}[Additivity in the universal cover]\label{rem:rotation-additivity}
% QC: this is the key property. rho(a~) lives in R (not S^1), so concatenation gives
%   exact addition, not just addition mod Z.
%   Contrast: if we just used the rotation angle mod 2pi, concatenation gives mod-Z addition.
%   The universal cover is essential for the pruning to work.
For paths $\tilde{a}$ and $\tilde{b}$ in $\widetilde{\operatorname{Sp}}(2)$
that compose (the endpoint of $\tilde{a}$ equals the basepoint of $\tilde{b}$):
\begin{equation}\label{eq:rotation-additivity}
  \rho(\tilde{a} \cdot \tilde{b}) = \rho(\tilde{a}) + \rho(\tilde{b}).
\end{equation}
The rotation number $\rho(\tilde{a}) \in \mathbb{R}$ is
the lift to $\mathbb{R}$, not the element of $S^1$:
it accumulates continuously and does not reduce modulo~$\mathbb{Z}$.
\end{remark}

\begin{definition}[Rotation increment at a breakpoint]\label{def:rotation-increment}
% QC: each directed 2-face j->l gives a transition from Reeb vector R_j to R_l.
%   The path f in Sp(2) tracking the linearized flow is constant along straight segments
%   (Remark~\ref{rem:rotation-polytope}) and jumps at breakpoints.
%   The jump at j->l is the transition matrix psi_{F_{jl}} from CH2021 §2.
%   By CH2021 Lemma~lem:transitionmatrix: psi is positive elliptic (trace in (-2,2), a2>0).
%   By CH2021 Corollary~cor:rotrange: rot_mod_Z(psi) in (0,1/2).
%   Since we work in the universal cover and start from rot=0 at t_1, each increment in (0,1/2).
% [TODO: JÖRN - state the exact formula for psi_{F_{jl}} from CH2021 (the explicit 2x2 matrix)
%   so the implementation can compute it. We need a reference to the precise formula.]
Let $F_j \to F_l$ be a directed $2$-face.
The \emph{rotation increment} $\Delta\rho_{jl} \in \mathbb{R}$
is the contribution to the rotation number at the breakpoint
where a Reeb trajectory transitions from velocity $R_j$ to $R_l$.
It is computed from the \emph{transition matrix} $\psi_{F_{jl}}$
of \cite{CH2021}, a $2 \times 2$ positive-elliptic symplectic matrix
associated with the $2$-face $F_{jl}$;
see \cite[Lem.~lem:transitionmatrix, Cor.~cor:rotrange]{CH2021}.
By those results, $\Delta\rho_{jl} \in (0, \tfrac{1}{2})$.
\end{definition}

\begin{remark}[Rotation on a polytope]\label{rem:rotation-polytope}
% QC: on a polytope, the linearized flow is constant (= identity) along straight Reeb segments,
%   because R_i is a constant vector (the flow map is a pure translation, Jacobian = I).
%   Rotation only changes at breakpoints. By additivity, total rotation = sum of increments.
%   This is why rho in the tube data is a single number: it's the sum over inner breakpoints.
On a polytope, the linearized flow of a pure Reeb trajectory
is the identity along straight segments
(since $R_i$ is constant, so the flow map $x \mapsto x + t R_i$ has Jacobian $I$).
By the additivity~\eqref{eq:rotation-additivity},
the total rotation of the tube $T(\sigma)$ from $t_1$ to $t_{k-1}$ is
\[
  \rho = \sum_{i=2}^{k-1} \Delta\rho_{\sigma(i-1),\sigma(i)},
\]
the sum of rotation increments at each inner breakpoint.
\end{remark}

\begin{lemma}[Rotation monotonicity]\label{lem:rotation-monotone}
% [TODO: JÖRN - this lemma claims each extension adds positive rotation.
%   The proof should follow from Delta_rho_{jl} > 0 (from the positive-elliptic character
%   of the transition matrix, CH2021 Corollary~cor:rotrange). But the increment Delta_rho
%   depends on the LIFT to the universal cover, not just the transition matrix value mod Z.
%   For the first increment, the lift starts at 0 so rho_1 = Delta_rho_{12} in (0,1/2).
%   For subsequent increments, rho accumulates. The claim is rho' > rho always.
%   Since Delta_rho_{jl} > 0, this is immediate. Verify that Delta_rho > 0 (not just mod-Z > 0)
%   for every extension step, not just the first.]
Each tube extension increases rotation:
writing $T_k = T(\sigma(1), \ldots, \sigma(k))$, if $T_k$
has rotation~$\rho$ then its extension $T_{k+1}$ has rotation
\[
  \rho' = \rho + \Delta\rho_{\sigma(k),\sigma(k+1)} > \rho.
\]
\end{lemma}

\begin{proof}
By Definition~\ref{def:rotation-increment},
$\Delta\rho_{\sigma(k),\sigma(k+1)} \in (0, \tfrac{1}{2})$,
so $\rho' > \rho$.
\end{proof}

% -------------------------------------------------------------------------
\subsection{Algorithm}\label{subsec:tube-algorithm-main}
% -------------------------------------------------------------------------

\begin{lemma}[Pruning: empty tube]\label{lem:prune-empty}
% QC: trivially correct. Empty Start or End means no trajectory of this type exists.
If $\operatorname{Start}(\sigma) = \emptyset$ or $\operatorname{End}(\sigma) = \emptyset$,
then $T(\sigma) = \emptyset$.
\end{lemma}

\begin{proof}
A trajectory $\gamma \in T(\sigma)$ satisfies
$\gamma(0) \in \operatorname{Start}(\sigma)$ and
$\gamma(t_{k-1}) \in \operatorname{End}(\sigma)$ by definition.
\end{proof}

\begin{lemma}[Pruning: action lower bound]\label{lem:prune-action}
% QC: min_{End} a is the minimum action over all trajectories in T(sigma).
%   Since a is affine on a convex polygon, the minimum is attained at a vertex.
%   If this minimum exceeds the best known candidate c*, no trajectory in T(sigma)
%   can improve the current best.
%   IMPORTANT: this compares the partial action (t_1 to t_{k-1}), not the full closed orbit action.
%   The pruning is still valid: the full closed orbit action = partial action + closing step time >= partial action.
%   [TODO: JÖRN - verify: can we prune on partial action < c* even before closing?
%   For tubes not yet closed, the partial action is a lower bound on the full action.
%   Yes: the full action = partial + closing time > partial. So if partial > c*, full > c* too.]
Let $c^* > 0$ be the action of the current best candidate orbit.
If $\min_{y \in \operatorname{End}} a(y) > c^*$,
then $T(\sigma)$ contains no minimum-action orbit.
\end{lemma}

\begin{proof}
Every $\gamma \in T(\sigma)$ has
$A(\gamma) \geq a(\gamma(t_{k-1})) \geq \min_{\operatorname{End}} a > c^*$.
Since the minimum-action orbit has action $\leq c^*$, it is not in $T(\sigma)$.
\end{proof}

\begin{lemma}[Pruning: rotation upper bound]\label{lem:prune-rotation}
% QC: the minimum-action orbit has rotation number strictly less than 2 (CH2021).
%   By rotation monotonicity (Lem~\ref{lem:rotation-monotone}), any further extensions
%   or closings only increase rho. So if rho >= 2 already, no closed orbit from this
%   tube can have rho < 2.
% [TODO: JÖRN - state and cite the exact result from CH2021 that the minimum-action
%   orbit has rho < 2. The relevant statement may be in CH2021 Corollary~cor:computecehz
%   or the paragraph at line~320 in s1. Verify the exact bound is < 2 (not ≤ 2 or ≤ 1).]
Let $\rho$ be the rotation number of tube $T(\sigma)$.
If $\rho \geq 2$, then $T(\sigma)$ contains no minimum-action orbit.
\end{lemma}

\begin{proof}
By \cite{CH2021}, the minimum-action combinatorial Reeb orbit
satisfies $\rho_{\mathrm{comb}} < 2$.
By Lemma~\ref{lem:rotation-monotone},
each extension and each closing step adds $\Delta\rho > 0$,
so the rotation of any closed orbit produced from $T(\sigma)$ exceeds $\rho \geq 2$.
\end{proof}

\begin{lemma}[Pruning: simple orbit]\label{lem:prune-simple}
% QC: minimum-action orbit is simple by \cite{HK2017} Theorem~1.1.
%   A simple orbit visits each facet on a contiguous time interval (by definition in HK2017).
%   This means no repeated facet index in sigma(1),...,sigma(k) (for the non-closing tube).
%   Exception: the closing step artificially appends sigma(1) and sigma(2) -- those
%   repetitions are expected and not a problem (they encode the periodicity condition).
If $\sigma(i) = \sigma(j)$ for some $1 \leq i < j \leq k$,
then $T(\sigma)$ contains no minimum-action orbit.
\end{lemma}

\begin{proof}
By \cite{HK2017}, the minimum-action generalized Reeb orbit is simple:
it visits each facet on a connected time interval and does not repeat facets.
A pure Reeb trajectory of type $\sigma$, when closed,
visits $F_{\sigma(i)}$ twice if $\sigma(i) = \sigma(j)$.
Hence no closed orbit from $T(\sigma)$ can be the minimum-action orbit.
\end{proof}

\begin{definition}[Closing a tube]\label{def:tube-close}
% QC: close T(sigma(1),...,sigma(k)) by appending sigma(k+1)=sigma(1) and sigma(k+2)=sigma(2).
%   Two extension steps: need F_{sigma(k)} -> F_{sigma(1)} and F_{sigma(1)} -> F_{sigma(2)}.
%   After closing: phi'' maps Start subset F_{sigma(1)} ∩ F_{sigma(2)} to F_{sigma(1)} ∩ F_{sigma(2)}.
%   A closed orbit in the tube satisfies gamma(t_1) = gamma(t_{k+1}), i.e. phi''(x) = x.
%   Action of the closed orbit through x: a''(phi''(x)) = a''(x) (since x is a fixed point).
% [TODO: JÖRN - the two extension steps append sigma(1) and sigma(2).
%   For the first extension: triple is (sigma(k-1), sigma(k), sigma(1)).
%   For the second: triple is (sigma(k), sigma(1), sigma(2)).
%   Need: F_{sigma(k)} -> F_{sigma(1)} and F_{sigma(1)} -> F_{sigma(2)} (directed edges exist).
%   These are the same directed edges as sigma(1)->sigma(2) and sigma(1) already in the tube.
%   Are these edges always valid? Yes, because sigma(1)->sigma(2) is the first edge of the tube.]
Given a tube $T(\sigma(1), \ldots, \sigma(k))$,
its \emph{closure} is obtained by performing two extension steps:
appending $\sigma(k+1) = \sigma(1)$ and then $\sigma(k+2) = \sigma(2)$.
The result is a tube $T(\sigma(1), \ldots, \sigma(k), \sigma(1), \sigma(2))$
with flow map $\varphi'' \colon \operatorname{Start} \to F_{\sigma(1)} \cap F_{\sigma(2)}$
and action function $a'' \colon \operatorname{End}'' \to \mathbb{R}$.
\end{definition}

\begin{lemma}[Closed orbits as fixed points]\label{lem:fixed-point}
% QC: a closed orbit in T(sigma) satisfies gamma(t_1) = gamma(t_{k+1}), i.e.
%   the flow map phi'' maps the starting point back to itself.
%   Fixed point: x in Start with phi''(x) = x (and x in End'' = phi''(Start)).
%   Action of the orbit through x = a''(x) = total time from t_1 to t_{k+1}.
%   If phi'' is a 2x2 affine map, fixed points are solutions to (phi'' - id)(x) = 0,
%   i.e. a linear system. Generically: unique fixed point (if phi''-id is invertible).
%   Degenerate: 0-dim, 1-dim, or 2-dim set of fixed points (e.g. Zoll polytopes).
The closed Reeb orbits in the closure of $T(\sigma)$ are in bijection with
fixed points of $\varphi'' \colon \operatorname{Start}' \to F_{\sigma(1)} \cap F_{\sigma(2)}$:
\[
  \gamma \text{ is a closed orbit in } T(\sigma)
  \;\Longleftrightarrow\;
  \varphi''(\gamma(0)) = \gamma(0)
  \;\text{ and }\;
  \gamma(0) \in \operatorname{Start}'.
\]
The action of the closed orbit through the fixed point $x$ is $a''(x)$.
\end{lemma}

\begin{proof}
A pure Reeb trajectory $\gamma \in T(\sigma(1), \ldots, \sigma(k))$ is closed
iff $\gamma(t_{k+1}) = \gamma(t_1)$.
After closure, $\gamma(t_{k+1}) = \varphi''(\gamma(0))$,
so the periodicity condition is $\varphi''(\gamma(0)) = \gamma(0)$.
The action is $a''(\varphi''(\gamma(0))) = a''(\gamma(0))$.
\end{proof}

\begin{algorithm}[Tube algorithm]\label{alg:tube}
% QC: Input must be a symplectic polytope (Definition~\ref{def:symplectic-polytope}).
%   Precomputation: O(F^3) step maps for all valid triples.
%   Search: DFS/BFS with pruning. Each node = tube data (Start, End, phi, a, rho).
%   Closing is attempted at each node, not just leaf nodes (early termination).
%   Correctness: Proposition~\ref{prop:tube-correctness}.
% Downstream: implement as recursive DFS over valid facet sequences.
%   Pruning lemmas as early-exit conditions.
%   Fixed-point computation: solve (phi'' - Id)(x) = 0 in local 2D coords.
Computes $c_{\mathrm{EHZ}}(K)$ for a symplectic polytope $K$
(assuming no Type~2 orbits; see Remark~\ref{rem:type-2}).

\smallskip\noindent
\emph{Input:} Facet data $(n_1, h_1), \ldots, (n_F, h_F)$.
\quad
\emph{Output:} $c_{\mathrm{EHZ}}(K)$.

\smallskip\noindent
\textbf{Precomputation.}
\begin{enumerate}
\item For each pair $(i,j)$: check if $F_i \cap F_j$ is a $2$-face
  and compute $\omega_0(n_i, n_j)$. Record directed edges $F_i \to F_j$ where $\omega_0(n_i, n_j) > 0$.
\item For each valid triple $(i,j,l)$ with $F_i \to F_j$ and $F_j \to F_l$:
  compute the step map $\Phi_{ijl}$~\eqref{eq:step-map},
  action increment $\Delta a_{ijl}$~\eqref{eq:action-increment},
  and rotation increment $\Delta\rho_{jl}$
  (Definition~\ref{def:rotation-increment}).
\end{enumerate}

\smallskip\noindent
\textbf{Search.}
Initialize $c^* = +\infty$.
For each directed edge $(i, j)$ with $F_i \to F_j$,
initialize tube $T(i, j)$ with data
$(\operatorname{Start}, \operatorname{End}, \varphi, a, \rho)$
as in Remark~\ref{rem:tube-init}.
Then recurse: given a tube $T(\sigma(1), \ldots, \sigma(k))$:

\begin{enumerate}
\item \textbf{Prune.} Discard the tube (return without recursing) if:
  \begin{itemize}
  \item $\operatorname{Start} = \emptyset$ or $\operatorname{End} = \emptyset$
    (Lemma~\ref{lem:prune-empty}), or
  \item $\min_{y \in \operatorname{End}} a(y) > c^*$
    (Lemma~\ref{lem:prune-action}), or
  \item $\rho \geq 2$
    (Lemma~\ref{lem:prune-rotation}).
  \end{itemize}

\item \textbf{Try closing.}
  Form the closure $T(\sigma(1), \ldots, \sigma(k), \sigma(1), \sigma(2))$
  (Definition~\ref{def:tube-close}).
  Find fixed points $x$ of $\varphi''$ in $\operatorname{Start}'$
  (Lemma~\ref{lem:fixed-point}).
  For each fixed point $x$, update $c^* \leftarrow \min(c^*, a''(x))$.

\item \textbf{Extend.}
  For each $l \notin \{\sigma(1), \ldots, \sigma(k)\}$
  with $F_{\sigma(k)} \to F_l$
  (Lemma~\ref{lem:prune-simple}):
  compute $T(\sigma(1), \ldots, \sigma(k), l)$
  (Definition~\ref{def:tube-extension}) and recurse.
\end{enumerate}

\smallskip\noindent
\textbf{Return} $(c^*)^{-1/2}$...
% [TODO: JÖRN - check: is c_EHZ = c* directly, or is there a factor?
%   In alg:ehz, the output is (max beta^T H beta)^{-1}. Here a(y) is the action (period)
%   of the Reeb orbit, i.e. A(gamma) = T. The EHZ capacity is the minimum action.
%   So c_EHZ = c* directly (no extra factor). But verify against Theorem~\ref{thm:main}.]
$c^*$.
\end{algorithm}

% -------------------------------------------------------------------------
\subsection{Correctness}\label{subsec:tube-correctness}
% -------------------------------------------------------------------------

\begin{proposition}[Tube algorithm correctness]\label{prop:tube-correctness}
% [TODO: JÖRN - this is a proof sketch. Full proof deferred to a separate section.
%   Key claims:
%   (1) HK2017 guarantees a minimum-action simple pure Reeb orbit exists.
%   (2) Such an orbit corresponds to some combinatorial type sigma and lies in T(sigma).
%   (3) The search tree is exhaustive: every T(sigma) is visited (up to pruning).
%   (4) Each pruning lemma is correct: it never discards a tube containing a min-action orbit.
%   Claim (4) needs verification for each lemma; claims (1-3) are standard.]
Let $K$ be a symplectic polytope with no Type~$2$ minimum-action orbit
(see Remark~\ref{rem:type-2} and \cite{CH2021}, Conjecture~conj:genericity).
Then Algorithm~\ref{alg:tube} computes $c_{\mathrm{EHZ}}(K)$ correctly.
\end{proposition}

\begin{proof}[Proof sketch]
By \cite{HK2017} (Theorem~\ref{thm:main}),
there exists a minimum-action generalized Reeb orbit $\gamma^*$
that is a simple pure Reeb orbit visiting each facet at most once.
Since $K$ has no Type~$2$ orbit by hypothesis,
every breakpoint of $\gamma^*$ lies in the interior of a $2$-face $F_{ij}$,
so $\gamma^*$ is of some type $\sigma^* = (\sigma^*(1), \ldots, \sigma^*(k^*))$.

The tube $T(\sigma^*)$ contains $\gamma^*$.
The search in Algorithm~\ref{alg:tube} eventually reaches $T(\sigma^*)$,
because it enumerates all sequences of directed edges
(depth-first, with pruning).
Each pruning lemma is sound:
Lemmas~\ref{lem:prune-empty} and~\ref{lem:prune-action} discard only empty
or suboptimal tubes;
Lemma~\ref{lem:prune-rotation} uses the bound
$\rho_{\mathrm{comb}}(\gamma^*) < 2$ from \cite{CH2021},
which $\gamma^*$'s tube does not violate;
Lemma~\ref{lem:prune-simple} uses simplicity of $\gamma^*$.
Hence $T(\sigma^*)$ survives all pruning steps.

When the closing step is applied to $T(\sigma^*)$,
Lemma~\ref{lem:fixed-point} yields $\gamma^*$ as a fixed point
with action $a^* = A(\gamma^*) = c_{\mathrm{EHZ}}(K)$.
The algorithm records $c^* = A(\gamma^*)$
and returns it.
\end{proof}

\begin{remark}[Type 2 orbits and non-generic polytopes]\label{rem:type-2}
% QC: CH2021 classifies combinatorial Reeb orbits into Type 1 (breakpoints at 2-face interiors),
%   Type 2 (breakpoints at 1-faces), Type 3 (contain a bad 1-face). See CH2021 §1.
%   Type 3 orbits do not contribute to c_EHZ (CH2021 Thm smoothtocomb(iii)).
%   Type 2 orbits: their breakpoints are at 1-faces, meaning F_{sigma(i)} ∩ F_{sigma(i+1)}
%   is a 1-face (not a 2-face). Such a pair (sigma(i), sigma(i+1)) would not appear as
%   a directed edge in Algorithm~\ref{alg:tube} (we only add edges for directed 2-faces).
%   Therefore, Algorithm~\ref{alg:tube} misses Type 2 orbits.
%   For generic symplectic polytopes (CH2021 Conjecture~conj:genericity), no Type 2 orbits exist.
%   The algorithm is therefore complete for generic polytopes.
%   Handling Type 2 orbits for non-generic polytopes: extend the search to include 1-face and
%   0-face transitions (at the cost of more complex step maps and rotation increments).
%   This is left as future work.
A \emph{Type~1 combinatorial Reeb orbit} (following \cite{CH2021}, \S 1)
is one whose breakpoints all lie in the interiors of $2$-faces.
A \emph{Type~2 orbit} has at least one breakpoint in the interior of a $1$-face.
(Type~3 orbits contain a ``bad'' $1$-face and do not contribute to $c_{\mathrm{EHZ}}$.)

Algorithm~\ref{alg:tube} searches only for Type~1 orbits:
each transition in the sequence $\sigma$ corresponds to a directed $2$-face $F_i \to F_j$,
and transitions through $1$-faces are not enumerated.
A Type~2 orbit would involve a transition $F_i \to F_l$ where
$F_i \cap F_l$ is a $1$-face (or $0$-face)
rather than a $2$-face.

By \cite[Conj.~conj:genericity]{CH2021},
generic symplectic polytopes have no Type~2 minimum-action orbit,
so Algorithm~\ref{alg:tube} is correct for generic polytopes.
Handling Type~2 orbits for non-generic polytopes is left as future work.
\end{remark}
